% =====================================================================
%  Lecture 24 -- Control Plane                            (English edition)
% =====================================================================
\documentclass[11pt,a4paper]{article}
\providecommand\ACNroot{../..}
\input{\ACNroot/templates/lecture-en}

\ACNset{number}{24}
\ACNset{title}{The 5G Control Plane}
\ACNset{shorttitle}{Control plane}
\ACNset{titleother}{صفحهٔ کنترل نسل پنجم}
\ACNset{part}{Part V -- Planes and outlook (Lectures 23--26)}
\ACNset{time}{6--8 hours (3 h reading, 3--5 h exercises)}
\ACNset{level}{Advanced}
\ACNset{prereq}{Lectures 4, 12, 18, 20 and 23}
\ACNset{gen}{5}

\begin{document}
\ACNcover

% ---------------------------------------------------------------------
\section{Learning objectives}
After studying this lecture you should be able to:
\begin{objectives}
  \item \textbf{Draw} the 5G control-plane protocol stack (RRC, NAS, NGAP, SBI) and \textbf{explain} which entities terminate each protocol.
  \item \textbf{Describe} the RRC states IDLE, INACTIVE and CONNECTED and their relation to the CM and RM states of the core.
  \item \textbf{Trace} a service request from CM-IDLE and \textbf{estimate} the control-plane latency of state transitions.
  \item \textbf{Compute} the signalling load on an AMF, the inactivity-timer trade-off and the optimal registration-area size.
  \item \textbf{Explain} how 5G protects its control plane against overload and signalling storms.
\end{objectives}

% ---------------------------------------------------------------------
\section{Prerequisite knowledge}
You need location management and paging (Lecture~4), the MME and its overload control (Lecture~12), the service-based architecture
(Lecture~18), the registration and PDU-session procedures (Lecture~20) and the user plane of Lecture~23. This lecture looks at the same
system from the signalling side: who talks to whom, in which state, and at what cost.

\paragraph{Key terms} The vocabulary of this lecture with its Persian equivalents:
\begin{keyterms}
Control plane & \fa{صفحهٔ کنترل} & Signalling that sets up, maintains and releases connectivity. \\
RRC & \fa{کنترل منابع رادیویی} & Radio Resource Control between UE and gNB. \\
NAS (MM / SM) & \fa{لایهٔ غیردسترسی} & Non-access stratum between UE and AMF (mobility) or SMF (sessions). \\
NGAP & \fa{پروتکل کاربرد NG} & Application protocol on N2 between gNB and AMF, over SCTP. \\
RRC\_INACTIVE & \fa{حالت غیرفعال RRC} & State in which the UE context is kept but the radio connection is suspended. \\
CM / RM state & \fa{حالت مدیریت اتصال / ثبت} & Core view of the UE: connected or idle, registered or deregistered. \\
Service request & \fa{درخواست خدمت} & NAS procedure that moves a UE from CM-IDLE to CM-CONNECTED. \\
Registration area & \fa{ناحیهٔ ثبت} & List of tracking areas in which an idle UE need not report its location. \\
RNA & \fa{ناحیهٔ اعلان مبتنی بر RAN} & Area in which an inactive UE moves without notifying the network. \\
Signalling storm & \fa{طوفان سیگنالینگ} & Burst of synchronised procedures that overloads control-plane nodes. \\
\end{keyterms}

% ---------------------------------------------------------------------
\section{Introduction and motivation}
The control plane does not carry a single user bit, yet it decides whether user bits can flow at all. Every time a smartphone wakes up to
fetch a notification, it runs a small protocol exchange across the radio, the gNB, the AMF, the SMF and the UPF. Multiplied by millions of
devices and dozens of transitions per hour, this signalling dominates the processing load of the core and a noticeable share of the radio
resources -- and its latency is felt by the user as the delay before the first packet.

5G reorganises the control plane in three ways: a new RRC state, RRC\_INACTIVE, that avoids most core signalling for short bursts; a clean
separation of mobility (AMF) and session (SMF) management; and service-based interfaces inside the core \cite{ts23501,ts38300}. This lecture
explains the resulting protocol stack and state machines and gives the models to dimension them.

% ---------------------------------------------------------------------
\section{Fundamental concepts}

\subsection{The control-plane stack}
Three protocol families make up the control plane (\cref{fig:cp}):
\begin{itemize}
  \item \textbf{RRC} between UE and gNB \cite{ts38331}: system information, connection setup, security activation on the radio,
        configuration of radio bearers, measurements and handover commands. It runs over PDCP, RLC and MAC on signalling radio bearers.
  \item \textbf{NAS} between UE and core \cite{ts24501}: NAS-MM (registration, service request, authentication, NAS security) terminates in
        the AMF; NAS-SM (PDU-session establishment and modification) is relayed by the AMF to the SMF. The gNB forwards NAS messages without
        reading them -- they are integrity-protected and ciphered end to end between UE and AMF.
  \item \textbf{NGAP} between gNB and AMF on N2 \cite{ts38413}: UE context management, PDU-session resource setup, paging and handover
        signalling, over SCTP. Inside the core, the AMF and SMF use service-based interfaces (HTTP/2 with JSON, Lecture~18).
\end{itemize}

\begin{acnfigure}
  \resizebox{\linewidth}{!}{\acnfig{lecture-24/cp-stack}}
  \caption{5G control-plane protocol stack: RRC on Uu, NGAP on N2, NAS end to end between UE and AMF/SMF, and SBI in the core.}
  \label{fig:cp}
\end{acnfigure}

\subsection{States of a UE}
The network keeps three coupled state machines for every UE \cite{ts23501,ts38331}:
\begin{itemize}
  \item \textbf{RM} (registration management) in the AMF: RM-DEREGISTERED or RM-REGISTERED.
  \item \textbf{CM} (connection management) in the AMF: CM-IDLE (no N2 connection for the UE; location known per registration area) or
        CM-CONNECTED (an N2 connection exists).
  \item \textbf{RRC} in the gNB and UE: RRC\_IDLE, RRC\_INACTIVE or RRC\_CONNECTED (\cref{fig:rrc}).
\end{itemize}
In RRC\_INACTIVE the UE behaves almost like an idle UE -- it camps on cells, monitors paging and moves freely inside a RAN-based notification
area (RNA) -- but the gNB that suspended it keeps the UE context and the N2/N3 connections to the core remain up, so the core sees the UE as
CM-CONNECTED. Returning to RRC\_CONNECTED takes an RRC Resume exchange with no core signalling.

\begin{acnfigure}
  \resizebox{0.92\linewidth}{!}{\acnfig{lecture-24/rrc-states}}
  \caption{NR RRC states and transitions; RRC\_INACTIVE belongs to CM-CONNECTED in the core.}
  \label{fig:rrc}
\end{acnfigure}

\subsection{Reachability}
Where the network can find the UE depends on its state: in CM-IDLE the AMF knows the registration area (a list of tracking areas) and must page
there; in RRC\_INACTIVE the anchor gNB knows the RNA and performs RAN paging itself; in RRC\_CONNECTED the serving cell is known. The UE reports
leaving its area by a mobility registration update (CM-IDLE) or an RNA update (RRC\_INACTIVE), and refreshes it periodically.

% ---------------------------------------------------------------------
\section{Detailed technical explanation}

\subsection{Service request}
When an idle UE has uplink data -- or is paged for downlink data -- it runs the service request (\cref{fig:sr}) \cite{ts23502}. After random
access the UE sets up an RRC connection and piggybacks the NAS Service Request on RRC Setup Complete; the gNB forwards it in an NGAP Initial
UE Message to the AMF selected by the 5G-S-TMSI. The AMF checks NAS integrity, asks the SMF of each PDU session to be activated to re-establish
its user plane (the SMF updates the UPF over N4), then sends an Initial Context Setup Request carrying the security key and the session
resources. The gNB activates AS security, reconfigures the radio bearers and returns its downlink tunnel endpoints, which the SMF installs in
the UPF.

\begin{acnfigure}
  \resizebox{\linewidth}{!}{\acnfig{lecture-24/service-request}}
  \caption{UE-triggered service request from CM-IDLE (simplified, after TS~23.502).}
  \label{fig:sr}
\end{acnfigure}

\subsection{Suspend and resume}
With RRC\_INACTIVE the gNB ends a connection by an RRC Release with suspend indication: the UE and the gNB keep the AS context and the
security configuration, the UE receives a resume identity and an RNA. When the UE has data again it sends an RRC Resume Request; if it is
in a cell of another gNB, that gNB retrieves the context from the anchor gNB over Xn and asks the AMF to switch the N2 path. In the common case
of resuming in the same area, no NAS or core signalling is needed at all. Small amounts of data can even be sent during the resume procedure
itself (small-data transmission, Release~17).

\subsection{Separation of mobility and session management}
In the EPC the MME handled both mobility and session signalling (Lecture~12). In 5G the AMF handles only mobility, access and NAS security,
while the SMF handles sessions. A UE has one AMF but may have several SMFs (one per PDU session or slice). This separation allows non-3GPP
access (Wi-Fi via N3IWF) to share the same AMF, keeps the AMF independent of session types, and lets slices have their own SMFs (Lecture~21).

\subsection{Overload and congestion control}
Control-plane nodes protect themselves at several levels \cite{ts23501,ts24501}: unified access control in the UE bars access categories
broadcast by the gNB before any message is sent; the AMF can ask gNBs to reject or reduce new RRC connections (NGAP overload start); NAS
congestion control returns a back-off timer (e.g.\ T3346) in reject messages so that UEs retry later; and the NRF and SCP spread load
over several AMF and SMF instances (Lecture~18).

% ---------------------------------------------------------------------
\section{Mathematical models}

\subsection{Signalling load}\label{sec:load}
\paragraph{Model} If $N$ UEs each perform procedure $k$ at rate $\lambda_k$ per hour and each procedure costs $m_k$ messages at the AMF, the
message rate in messages per second is
\begin{equation}\label{eq:load}
  \Lambda = \frac{N}{3600} \sum_{k} \lambda_k\, m_k .
\end{equation}
\paragraph{Example} $N = 2\times10^{6}$ with busy-hour rates per UE: 20 service requests (4 messages each), 2 mobility registrations (6), 3
handovers (6), 0.1 initial registrations (20) and 0.5 PDU-session establishments (8): $\sum_k \lambda_k m_k = 116$ messages per UE per hour and
$\Lambda = 64\,444$ messages/s. If RRC\_INACTIVE turns 80\,\% of the idle-to-connected transitions into resumes without core signalling, the sum
falls to 52 and $\Lambda = 28\,889$ messages/s -- a reduction of 55\,\%.

\subsection{Inactivity timer}\label{sec:timer}
\paragraph{Model} A UE generates data bursts as a Poisson process of rate $\lambda$. After each burst the gNB keeps the connection for an inactivity
time $T$ before releasing (or suspending) it. A transition to idle happens after a burst only if the next burst comes later than $T$, so the
transition rate and the mean connected ``tail'' per burst are
\begin{equation}\label{eq:timer}
  r(T) = \lambda\, e^{-\lambda T}, \qquad \bar{t}_{\mathrm{tail}}(T) = \frac{1 - e^{-\lambda T}}{\lambda} .
\end{equation}
The fraction of time spent in the tail is $\lambda\,\bar{t}_{\mathrm{tail}} = 1 - e^{-\lambda T}$.
\paragraph{Example} One burst per minute ($\lambda = 1/60$\,s$^{-1}$): $T = 10$\,s gives 50.8 transitions per hour and 15.4\,\% of the time in the
tail; $T = 30$\,s gives 36.4 transitions per hour and 39.3\,\% (\cref{fig:timer}). A longer timer saves signalling but costs UE energy and radio
resources; RRC\_INACTIVE makes short timers affordable because each transition is cheap.

\begin{acnfigure}
  \acnfig{lecture-24/inactivity-timer}
  \caption{Inactivity-timer trade-off for one burst per minute (\cref{eq:timer}).}
  \label{fig:timer}
\end{acnfigure}

\subsection{Registration-area size}\label{sec:ra}
\paragraph{Model} With the fluid-flow mobility model (Lecture~4), a UE moving at mean speed $v$ in an area of perimeter $P$ and size $A$ crosses
its border at rate $vP/(\pi A)$. For a square area of side $a$ this is $R_u = 4v/(\pi a)$. If each update costs $w_u$ and each paging event
(rate $\lambda_p$) costs $w_p$ in each of the $n = a^2/A_c$ cells of size $A_c$, the cost per UE and hour is
\begin{equation}\label{eq:racost}
  C(a) = w_u \frac{4v}{\pi a} + w_p \lambda_p \frac{a^2}{A_c} ,
\end{equation}
minimised by
\begin{equation}\label{eq:raopt}
  a^{*} = \left( \frac{2 w_u v A_c}{\pi w_p \lambda_p} \right)^{1/3} .
\end{equation}
At the optimum the update cost is exactly twice the paging cost.
\paragraph{Example} $v = 20$\,km/h, $A_c = 1$\,km$^2$, $\lambda_p = 5$ per hour, $w_u = 10$, $w_p = 1$: $a^{*} = 2.94$\,km (about 8.7 cells per area),
with update cost 86.6 and paging cost 43.3 per UE and hour (\cref{fig:ra}).

\begin{acnfigure}
  \acnfig{lecture-24/ta-cost}
  \caption{Registration-area size trade-off (\cref{eq:racost,eq:raopt}).}
  \label{fig:ra}
\end{acnfigure}

\subsection{Control-plane latency}\label{sec:cplat}
\paragraph{Model} The time from a battery-efficient state to the start of data transfer is the sum of the steps of the transition:
\begin{equation}\label{eq:cplat}
  T_{\mathrm{CP}} = T_{\mathrm{RA}} + T_{\mathrm{RRC}} + T_{\mathrm{core}} + T_{\mathrm{ctx}} .
\end{equation}
TR~38.913 sets a target of 10\,ms \cite{tr38913}; ITU-R M.2410 requires 20\,ms and encourages 10\,ms \cite{itu2017m2410}.
\paragraph{Example} With illustrative values for 30\,kHz subcarrier spacing -- random access 2\,ms, RRC exchange 4\,ms, core signalling (N2, AMF,
SMF and N4) 6\,ms, context setup with AS security and reconfiguration 4\,ms -- the service request from RRC\_IDLE takes 16\,ms. A resume from
RRC\_INACTIVE needs only random access and the RRC exchange: 6\,ms. Only the inactive state meets the 10\,ms target.

% ---------------------------------------------------------------------
\section{Network architecture and diagrams}
The control plane spans the UE, the gNB (CU-CP after a split, Lecture~23), the AMF set serving a region, and the SMFs, PCF, UDM and AUSF
reached through the SBA. AMFs are deployed as sets of interchangeable instances sharing UE contexts in a storage function, so that a gNB can
reach any instance of the set (a GUAMI identifies the set and pointer). \Cref{fig:cp,fig:rrc,fig:sr} show the protocol view, the state view and
the procedure view of this architecture.

% ---------------------------------------------------------------------
\section{Step-by-step procedures}

\begin{procedurebox}{Mobile-terminated data for an idle UE}\label{proc:mt}
\begin{steps}
  \item A downlink packet arrives at the UPF for a session without an N3 tunnel; the UPF buffers it (or forwards it to the SMF) and notifies the SMF.
  \item The SMF asks the AMF to make the UE reachable (Namf\_Communication\_N1N2MessageTransfer).
  \item The AMF pages the UE in all cells of its registration area through the gNBs (NGAP Paging).
  \item The UE answers with a service request (\cref{fig:sr}); the user plane is re-established and the buffered data are delivered.
\end{steps}
\end{procedurebox}

\begin{procedurebox}{Resuming from RRC\_INACTIVE}\label{proc:resume}
\begin{steps}
  \item The UE (with uplink data, or after RAN paging) performs random access and sends an RRC Resume Request with its resume identity.
  \item If the context is held by another gNB, the new gNB retrieves it over Xn and later asks the AMF for an N2 path switch.
  \item The gNB answers with RRC Resume; security is restored from the stored context with fresh keys.
  \item The UE confirms with RRC Resume Complete and data flow on the existing N3 tunnels.
\end{steps}
\end{procedurebox}

% ---------------------------------------------------------------------
\section{Worked examples}

\begin{examplebox}{Dimensioning an AMF set}\label{ex:load}
\textbf{Problem.} For the load of \cref{sec:load}, how many AMF instances are needed if each processes $\mu = 20\,000$ messages/s and should run at no
more than 70\,\% utilisation? Compare with and without RRC\_INACTIVE.

\textbf{Solution.} Usable capacity per instance $0.7\mu = 14\,000$ messages/s. Without INACTIVE: $64\,444/14\,000 = 4.6$, so 5 instances. With INACTIVE:
$28\,889/14\,000 = 2.06$, so 3 instances. Signalling efficiency in the RAN translates directly into fewer core instances.
\end{examplebox}

\begin{examplebox}{Choosing the inactivity timer}\label{ex:timer}
\textbf{Problem.} With one burst per minute, what is the longest timer that keeps the UE in the tail for at most 20\,\% of the time, and how many
transitions per hour result?

\textbf{Solution.} $1 - e^{-T/60} \le 0.2 \Rightarrow T \le -60\ln 0.8 = 13.4$\,s. Transitions: $60\,e^{-13.4/60} = 60\times0.8 = 48$ per hour.
\end{examplebox}

\begin{examplebox}{Registration area for fast users}\label{ex:ra}
\textbf{Problem.} Repeat the example of \cref{sec:ra} for vehicular users with $v = 60$\,km/h.

\textbf{Solution.} $a^{*3} = 2\times10\times60\times1/(\pi\times5) = 76.4$, $a^{*} = 4.24$\,km (about 18 cells). Update cost $10\times4\times60/(\pi\times4.24) = 180$,
paging cost $5\times4.24^2 = 90$, total 270 per UE and hour. Faster users call for larger areas, which is why operators give TA lists along motorways
and railways more cells.
\end{examplebox}

% ---------------------------------------------------------------------
\section{Technology comparison}
\begin{table}[htbp]
\centering\small
\caption{Control plane in LTE/EPC and in NR/5GC.}\label{tab:cmp}
\rowcolors{2}{acnmist}{white}
\begin{tabularx}{\linewidth}{@{}l Y Y@{}}
\toprule
\rowcolor{white} & LTE / EPC & NR / 5GC \\
\midrule
RRC states         & IDLE, CONNECTED (suspend/resume for IoT only) & IDLE, INACTIVE, CONNECTED \\
Mobility and sessions & MME (both), S-GW/P-GW control & AMF (mobility), SMF (sessions) \\
Core interfaces    & Diameter, GTPv2-C (point to point) & SBI over HTTP/2, NRF discovery \\
RAN--core protocol & S1AP over SCTP                  & NGAP over SCTP \\
Location areas     & TA lists, paging by MME          & TA lists (CM-IDLE) and RNA with RAN paging (INACTIVE) \\
Access control     & ACB, EAB, back-off timers       & Unified access control, back-off timers \\
Non-3GPP access    & Separate ePDG path                & Same AMF via N3IWF \\
\bottomrule
\end{tabularx}
\end{table}

\paragraph{Signalling is the price of efficiency} Each idle period saves energy and radio resources but must be paid for by a transition. The
5G control plane lowers that price (INACTIVE, fewer core messages, faster processing), which lets operators release connections earlier and
serve many more devices per unit of core capacity.

% ---------------------------------------------------------------------
\section{Key takeaways}
\begin{takeawaybox}
\begin{itemize}
  \item RRC (UE--gNB), NAS-MM (UE--AMF), NAS-SM (UE--SMF via AMF), NGAP (gNB--AMF over SCTP) and SBI in the core.
  \item RRC\_INACTIVE keeps the context in the RAN and the core in CM-CONNECTED: resume needs no core signalling.
  \item Load $\Lambda = (N/3600)\sum_k \lambda_k m_k$; idle-to-connected transitions dominate it.
  \item Inactivity timer: transitions $\lambda e^{-\lambda T}$ versus tail fraction $1 - e^{-\lambda T}$.
  \item Registration area: $a^{*} = (2w_u v A_c/(\pi w_p\lambda_p))^{1/3}$; updates cost twice the paging at the optimum.
\end{itemize}
\end{takeawaybox}

% ---------------------------------------------------------------------
\section{Exercises}

\begin{exercisebox}{1}{Where is the UE?}\label{exr:states}
For each RRC state, give the CM state seen by the AMF and the granularity at which the network knows the UE's location.
\end{exercisebox}

\begin{exercisebox}{1}{Transparent NAS}\label{exr:nas}
Why is NAS transparent to the gNB? Give a security reason and an architectural reason.
\end{exercisebox}

\begin{exercisebox}{2}{AMF load}\label{exr:load}
An AMF set serves $5\times10^{5}$ UEs with busy-hour rates per UE of 30 service requests (4 messages each), 1 mobility registration (6) and 4 handovers (6).
Compute the message rate.
\end{exercisebox}

\begin{exercisebox}{2}{Chatty applications}\label{exr:timer}
A messaging application produces one burst every 20\,s on average. Compute the transitions per hour and the tail fraction for $T = 5$\,s and $T = 10$\,s.
\end{exercisebox}

\begin{exercisebox}{3}{Optimal area}\label{exr:ra}
Derive \cref{eq:raopt} from \cref{eq:racost}, show that the update cost is twice the paging cost at the optimum, and compute $a^{*}$ for pedestrians
($v = 5$\,km/h) with the other values of \cref{sec:ra}.
\end{exercisebox}

\begin{exercisebox}{3}{Signalling storm}\label{exr:storm}
After a core outage, $10^{6}$ UEs must register again, each registration costing 20 messages at an AMF set of 5 instances with $\mu = 20\,000$ messages/s.
If the UEs spread their attempts uniformly over a window $W$, what is the smallest $W$ for 100\,\% and for 70\,\% utilisation? Which mechanisms enforce such a spread?
\end{exercisebox}

\subsection*{Solutions}
\addcontentsline{toc}{subsection}{Solutions}

\begin{solutionbox}{\ref{exr:states}}
RRC\_IDLE: CM-IDLE, registration area (list of TAs). RRC\_INACTIVE: CM-CONNECTED, RNA known to the anchor gNB (the AMF knows the gNB). RRC\_CONNECTED: CM-CONNECTED,
serving cell.
\end{solutionbox}

\begin{solutionbox}{\ref{exr:nas}}
Security: NAS messages are protected end to end with keys shared only by UE and AMF, so a compromised or untrusted RAN node (including non-3GPP access) cannot read
or alter them. Architecture: the RAN is independent of core procedures; new NAS features need no RAN change, and the same NAS runs over 3GPP and non-3GPP access.
\end{solutionbox}

\begin{solutionbox}{\ref{exr:load}}
$\sum_k \lambda_k m_k = 120 + 6 + 24 = 150$ per UE and hour; $\Lambda = 5\times10^5\times150/3600 = 20\,833$ messages/s.
\end{solutionbox}

\begin{solutionbox}{\ref{exr:timer}}
$\lambda = 1/20$\,s$^{-1}$, 180 bursts per hour. $T = 5$\,s: $180\,e^{-0.25} = 140.2$ transitions per hour, tail $1 - e^{-0.25} = 22.1\,\%$. $T = 10$\,s: $180\,e^{-0.5} = 109.2$,
tail 39.3\,\%. Frequent small bursts are the case for which RRC\_INACTIVE and small-data transmission were designed.
\end{solutionbox}

\begin{solutionbox}{\ref{exr:ra}}
$C'(a) = -4w_uv/(\pi a^2) + 2w_p\lambda_p a/A_c = 0 \Rightarrow a^3 = 2w_uvA_c/(\pi w_p\lambda_p)$. At the optimum $w_p\lambda_p a^2/A_c = 2w_uv/(\pi a)$, i.e.\ paging cost $= 2w_uv/(\pi a)$,
while update cost $= 4w_uv/(\pi a)$: twice as much. Pedestrians: $a^{*3} = 2\times10\times5/(5\pi) = 6.37$, $a^{*} = 1.85$\,km.
\end{solutionbox}

\begin{solutionbox}{\ref{exr:storm}}
Total $2\times10^{7}$ messages; capacity $10^{5}$ messages/s. 100\,\%: $W \ge 200$\,s; 70\,\%: $W \ge 2\times10^{7}/(7\times10^{4}) = 286$\,s. Mechanisms: randomised back-off timers in NAS
rejects (e.g.\ T3346), unified access control barring broadcast by the gNBs, AMF overload indications to the RAN, and prioritisation of emergency and high-priority
access.
\end{solutionbox}

% ---------------------------------------------------------------------
\section{Further reading}
The RRC protocol and states are specified in TS~38.331 \cite{ts38331} and summarised in TS~38.300 \cite{ts38300}; NAS in TS~24.501 \cite{ts24501}; NGAP in TS~38.413
\cite{ts38413}; the CM and RM states in TS~23.501 \cite{ts23501} and the procedures in TS~23.502 \cite{ts23502}. The latency targets come from TR~38.913 \cite{tr38913} and
ITU-R M.2410 \cite{itu2017m2410}. Rommer et al.\ \cite{rommer2020core} explain the 5G core signalling; Dahlman et al.\ \cite{dahlman2020nr} the RRC states and small-data
transmission.

\printbibliography[heading=subbibliography,title={\ACNstrRefs}]
\end{document}
